% =====================================================================
% ROZDZIAŁ 2 — Przegląd literatury
% =====================================================================
\chapter{Przegląd literatury}
\label{rozd:przeglad}

Przegląd grupuje literaturę wokół miejsca, w którym rozprawa wnosi
wkład. Sekcja~\ref{sec:p-ogon} opisuje rozjazd między detekcją phishingu
masowego a ogonem ukierunkowanym; sekcja~\ref{sec:p-behawioralna} ---
metody behawioralne i nagłówkowe, najbliższe technicznej rodzinie
detektorów tej rozprawy; sekcja~\ref{sec:p-osint} --- wiedzę
atakującego i narażenie ludzkie, czyli stronę ataku
(spear phishingu jako naruszenia grafu wiedzy); sekcje
\ref{sec:p-multiplex} i \ref{sec:p-gnn} --- formalizmy i modele
grafowe; sekcja~\ref{sec:p-rygor} --- rygor ewaluacji w bezpieczeństwie
uczenia maszynowego, do którego rozdział formalny dodaje twierdzenia.

\section{Detekcja phishingu: ogon masowy i ogon ukierunkowany}
\label{sec:p-ogon}

Nowoczesne filtry pocztowe klasyfikują kampanie masowe z dokładnością
rzędu 98\% na typowych benchmarkach \citep{birthriya2025detection,
tan2020graph}, łącząc cechy treści, URL-i i reputacji. Dokładność ta
jest jednak zdominowana przez ogon masowy i nie przenosi się na rzadki
ogon ukierunkowany: phishing spear i BEC są z definicji spersonalizowane
\citep{birthriya2025detection}, a ich treść czyta się jako normalna.
\citet{cidon2019high} wykazali, że BEC bez ładunku (payload-free) ---
wiadomość bez linków i załączników --- stanowi w produkcji skuteczną
i kosztowną klasę ataków, omijającą filtry treści;
\citet{ho2019detecting} scharakteryzowali lateral phishing na
113~milionach wiadomości, ustalając jednocześnie operacyjny próg
praktyczny: mniej niż 4 fałszywe alarmy na milion wiadomości
($4\cdot10^{-6}$). Straty BEC szacuje IC3 na $\sim$2{,}9 mld USD
w 2023 r. \citep{fbi2024ic3}. Wniosek wspólny tej linii badań brzmi:
dla ogona ukierunkowanego sygnał treściowy jest z definicji słaby,
bo atak respektuje wszystko, czego ofiara oczekuje
\citep{butavicius2022why,sturman2025security,dewan2014analyzing},
a koszt personalizacji obniżają językowe modele
\citep{bethany2024evaluating,brown2025web}.

\section{Detekcja behawioralna i nagłówkowa poczty}
\label{sec:p-behawioralna}

Metody skuteczne dla ogona ukierunkowanego odchodzą od treści.
Profile behawioralne per nadawca --- COMPA \citep{egele2013compa}
oraz statystyki „that ain't you" \citep{stringhini2015thataint} ---
flagują wysyłki odbiegające od własnego profilu aktywności;
\citet{gascon2018reading} modelują rytm pisania,
\citet{duman2016emailprofiler} łączą nagłówki ze stylometrią,
a \citet{luo2024characterizing} analizują infrastrukturę przekazań
(relay). W rodzinie logów systemowych Hopper \citep{ho2021hopper}
wykrywa lateral movement w grafach logowań, czytając ścieżkę
przyczynową do rzadkiego celu --- ideę, którą artykuł~A przenosi
na graf poczty jako baseline path-based. Wszystkie te detektory
czytają jednak \emph{stronę nadawcy}: to, co atakujący kontroluje
i świadomie rozkłada w czasie. Wkład rozprawy względem tej linii
jest dwojaki: (i)~sygnał lokowany jest po \emph{stronie odbiorczej}
--- w falach wśród sąsiadów ofiary, których atakujący nie może
ukryć bez rezygnacji z fan-outu (rozdz.~\ref{rozd:kaskada});
(ii)~rytm nadawcy, serce detekcji behawioralnej, okazuje się
jednocześnie najłatwiejszym miejscem do nieumyślnego wycieku
w~ewaluacji (pułapka wycieku rytmu, rozdz.~\ref{rozd:multipleks};
formalizacja w tw.~\ref{thm:lookahead}).

\section{Strona ataku: OSINT, personalizacja i narażenie ludzkie}
\label{sec:p-osint}

Spear phishing czerpie wiarygodność z otwartych źródeł (OSINT);
automatyczne zbieranie śladów badano od wczesnych prac
\citep{edwards2017panning}. \citet{dewan2014analyzing,dewan2017facebook}
znaleźli jednak, że surowe cechy społecznościowe i OSINT wnosiły
mało ponad treść i URL-e --- ostrzeżenie, które rozprawa bierze
poważnie: OSINT działa jako \emph{konstruktor relacji} (wspólna
lista czy afiliacja = krawędź warstwy), a nie płaska cecha
(binarna flaga). Modele językowe radykalnie obniżyły koszt
personalizacji: automatyczne kampanie dorównują ekspertom
w badaniach z udziałem ludzi \citep{bethany2024evaluating},
a realistyczne ślady OSINT można generować syntetycznie
\citep{brown2025web}. Razem z tym, że treść może być wygenerowana
idealnie, wynika imperatyw: obrona musi eksploatować \emph{relacje},
nie styl.

Narażenie ludzkie zależy od osobowości i kontekstu
\citep{halevi2015spearphishing}, rośnie pod presją czasu
i sygnałami perswazji \citep{butavicius2022why}, a świadomość
i styl decyzyjny moderują je tylko częściowo \citep{sturman2025security}.
Te prace mierzą stronę \emph{ludzką}; rozprawa mierzy wielkość
komplementarną i strukturalną --- narażenie, które \emph{detektor
zostawia} dla danej roli i kanału ataku (def. susceptibility,
rozdz.~\ref{rozd:multipleks}).

\section{Sieci multiplex i niespójność międzywarstwowa}
\label{sec:p-multiplex}

Formalizmem scalającym wiele relacji na jednym zbiorze węzłów jest
sieć multiplex \citep{kivela2014multilayer,dedomenico2013mathematical,
magnani2013combinatorial,jeude2019multilayer}. W modelach uczonych
heterogeniczne GNN uczą się fuzji typów relacji: HAN
\citep{wang2019heterogeneous}, HGT \citep{hu2020heterogeneous},
MAGNN \citep{fu2020magnn}, przegląd reprezentacji sieci
heterogenicznych \citep{christophides2021overview,cen2019representation}.
Niespójność między perspektywami jako sygnał anomalii znana jest
w wykrywaniu oszustw \citep{liu2020alleviating}; ataki na strukturę
grafu \citep{zugner2018adversarial,zugner2019adversarial,
dai2022towards}, ich certyfikowane granice odporności
\citep{zugner2021certified} i obrony \citep{zhang2020gnnguard,jin2020graph,
cai2021structural,wang2023label} mają znaczenie, bo atakujący może
fabrykować krawędzie --- a w grafach wiedzy robi to na dwóch
poziomach: trucie osadzeń i korpusów GraphRAG
\citep{zhang2019poisoning,you2023mass,zou2024poisonedrag,
wen2025fewwords,kepo2026} oraz wandalizm na publicznych grafach
wiedzy \citep{heindorf2017wsdmcup,sarabadani2017tools,
heindorf2019debiasing}. Co rozdziały formalny (stw.
o fabrykacji, \ref{prop:fab}) i empiryczne B i C (sweep fabrykacji)
traktują jako obywatela pierwszej kategorii modelu zagrożeń.
Według wiedzy autora niespójność międzywarstwowa nie została
dotąd operacjonalizowana jako graf wiedzy domenowej poczty
korporacyjnej dla bezpieczeństwa; tę lukę zamyka rozdział
\ref{rozd:multipleks}.

\section{Grafowe modele uczone w czasie}
\label{sec:p-gnn}

Tematyczne sieci grafowe uczą reprezentacje na grafie statycznym;
ich temporalny odpowiednik utrzymuje pamięć węzłów aktualizowaną
kolejnymi interakcjami --- klasę reprezentuje TGN
\citep{rossi2020temporal}. W rozprawie architektura temporalna
z pamięcią węzłów (TGN-lite) jest użyta jako \emph{instrument
pomiarowy}: różnica między nią a modelem statycznym (i spadek
tej różnicy pod tasowaniem czasu) operacjonalizuje tezę, że
sygnał kaskadowy jest dynamiczny, nie topologiczny. Graf
komunikacji jako obiekt monitoringu operacyjnego jest tej samej
natury co grafy przepływów, alertów i proweniencji stosowane
w detekcji anomalii i intruzji \citep{nayeri2026alert,weng2024rtapt}.

\section{Rygor ewaluacji detektorów bezpieczeństwa}
\label{sec:p-rygor}

Wyniki detektorów ML w bezpieczeństwie łatwo ulegają zawyżeniu: wycieki cech,
nieuczciwe podziały i nagłówkowy AUC \citep{arp2022dos};
świadome czasu ewaluacje pokazują, jak dylatacja czasowa
zniekształca wyniki (TESSERACT, \citep{pendlebury2019tesseract}).
Ta linia dostarcza protokołów; rozprawa idzie krok dalej ---
aksjomatyzuje warunki w postaci twierdzeń o obciążeniu
(tw. o look-ahead \ref{thm:lookahead}, stw. o podziale po
ofiarach \ref{prop:split}, lemat o sondzie \ref{lem:probe},
tw. o usunięciu sygnatury \ref{thm:injection}, stw. o
kontroli tasowania \ref{prop:shuffle}, stw. o fabrykacji
\ref{prop:fab}) i pokazuje, że pułapka wycieku rytmu
(rozdz.~\ref{rozd:multipleks}) jest \emph{instancją}
dokładnie tego mechanizmu, w pełni policzalną na realnym
korporusie Enron \citep{klimt2004enron}.

\section{Podsumowanie luki badawczej}
\label{sec:p-luka}

Z przeglądu wyłaniają się trzy luki, które rozprawa zamyka.
Po pierwsze, metody dla ogona ukierunkowanego czytają stronę
nadawcy (wolumetria, rytm, ścieżka) --- brakowało pomiaru sygnału
po stronie odbiorczej fali kaskady, niewidocznego dla atakującego
bez rezygnacji z fan-outu. Po drugie, wiedza organizacji
była modelowana płasko (cechy), nie relationalnie (warstwy grafu
wiedzy) --- brakowało operacjonalizacji niespójności międzywarstwowej
dla impersonacji zewnętrznej. Po trzecie, rygor ewaluacji był
protokołowy, nie twierdzeniowy --- brakowało formalnych warunków,
pod jakimi empiryczny AUC detektora strumieniowego mówi cokolwiek
o produkcji, oraz pomiaru, jak duży jest zysk z ich naruszenia
na realnym korpusie.
